\documentclass[12pt]{article}
\usepackage[T1]{fontenc}
\usepackage[utf8]{inputenc}
\usepackage{lmodern}
\usepackage{textcomp}
\usepackage{enumitem}
\usepackage{geometry}
\usepackage{graphicx}
\usepackage{amsmath, amssymb}
\geometry{margin=1in}
\begin{document}
\begin{center}
Religious Studies - K-12 Level Assessment 16 \
Total Marks: 40 \
Answer all questions.
\end{center}

\section*{Multiple Choice (10 marks)}
\begin{enumerate}[label=\arabic*.]
\item The term xin refers to which of the following?
\begin{enumerate}[label=(\alph*)]
    \item Piety
    \item Non-action
    \item Heart-mind
    \item World
\end{enumerate}
\item What is the English title for the divination text also known as the Yijing?
\begin{enumerate}[label=(\alph*)]
    \item The Classic of Documents
    \item The Lotus Sutra
    \item The Flower Garland Sutra
    \item The Classic of Changes
\end{enumerate}
\item When was the current Dalai Lama born?
\begin{enumerate}[label=(\alph*)]
    \item 1935
    \item 2000
    \item 1965
    \item 1975
\end{enumerate}
\item Which term refers to the five main areas into which early Christianity was divided?
\begin{enumerate}[label=(\alph*)]
    \item Presbyteries
    \item Sees
    \item Churches
    \item Synods
\end{enumerate}
\item Which Japanese government promoted a kind of national cult based on the emperor and his associations with kami?
\begin{enumerate}[label=(\alph*)]
    \item Honen
    \item Tanaka
    \item Tokugawa
    \item Meiji
\end{enumerate}
\end{enumerate}

\section*{True / False (10 marks)}
\textit{Answer True or False}
\begin{enumerate}[label=\arabic*.]
\setcounter{enumi}{5}
\item True or False: Religious traditions in China and Korea primarily emphasize the pursuit of power and influence over others.
\item True or False: Ordinary folk in the Han Dynasty commonly sought assistance from Laozi during times of drought.
\item True or False: In the Zoroastrian tradition, the term Gathas refers to a set of ethical laws governing behavior.
\item True or False: The Hebrew word mashiach means 'prophet' and signifies a religious teacher in Jewish culture.
\item True or False: Gandhi's title 'Mahatma' translates to 'Great soul' in English.
\end{enumerate}

\section*{Long Form Response (20 marks)}
\textit{Answer each question with a clear and complete explanation.}
\begin{enumerate}[label=\arabic*.]
\setcounter{enumi}{10}
\item Discuss the significance of Hwanung's transformation into a human in the context of the Korean foundation myth and its implications for understanding Korean cultural beliefs.
\item Discuss the significance of Guru Nanak starting to preach the message of the divine Name at the age of 30 and how it influenced the development of Sikhism.
\end{enumerate}

\vfill
\noindent\textit{End of Paper}
\end{document}